\section{Direct posterior bounds and interpolation between Gaussian scales}
\label{app:scale-interpolation}

The finite message used by the entropy-metric reconstruction has a
quadratic dependence on the active rank. Its accumulated cost need not
be converted into a dense small-width code. The actual conditioning
posterior can be treated directly whenever that message cost is larger
than its Gaussian relative-entropy bound. We give a single statistical
event for this argument and for the scale-dependent reconstruction.
All reference laws below are fixed before the training sample is drawn.
This section prices the comparison class uniformly. The later
Gaussian-weighted construction improves its all-scale rate and removes
the intermediate exponent restriction appearing in this preliminary
calculation.

\begin{lemma}[A relative-entropy risk inequality]
\label{lem:si-risk}
Let $\Pi$ be a probability law on a measurable family of predictors,
and let $\ell_h(x)\in[0,1]$. Write
$R(h)=\mathbb E_X\ell_h(X)$ and
$\widehat R_S(h)=n^{-1}\sum_{i=1}^n\ell_h(X_i)$.
With sample probability at least $1-\delta$, simultaneously for every
probability law $\Gamma\ll\Pi$,
\begin{equation}\label{eq:si-risk}
 \mathbb E_\Gamma R
 \le 2\mathbb E_\Gamma\widehat R_S
       +\frac{2}{n}
          \left[D_{\rm KL}(\Gamma\Vert\Pi)+\log(1/\delta)\right].
\end{equation}
\end{lemma}
\begin{proof}
Put $a=1-\exp(-1)>1/2$. Convexity gives
$e^{-u}\le1-au$ for $0\le u\le1$. For each fixed predictor,
independence of the sample therefore gives
\[
 \mathbb E_S\exp\{n[aR(h)-\widehat R_S(h)]\}
 \le e^{naR(h)}[1-aR(h)]^n\le1.
\]
Fubini and Markov's inequality imply, with probability at least
$1-\delta$,
\[
 \int\exp\{n[aR(h)-\widehat R_S(h)]\}\,d\Pi(h)
 \le1/\delta.
\]
On this event, Jensen's inequality applied under $\Gamma$ gives
\[
 n\left[a\mathbb E_\Gamma R-
                   \mathbb E_\Gamma\widehat R_S\right]
 \le D_{\rm KL}(\Gamma\Vert\Pi)+\log(1/\delta).
\]
Indeed, the general inequality used here is
$\int F\,d\Gamma-D_{\rm KL}(\Gamma\Vert\Pi)
 \le\log\int e^F\,d\Pi$; it follows by applying concavity of the
logarithm to $e^F(d\Pi/d\Gamma)$. Division by $a>1/2$ proves
\eqref{eq:si-risk}. Infinite relative entropy makes the assertion
vacuous and causes no exception.
\end{proof}

The target cylinder gives, for the actual posterior,
\begin{equation}\label{eq:si-raw-density}
 \frac{dQ_{S,N,t,c}}{d\mu_{N,t}}
 \le\exp(NK/t^2).
\end{equation}
Here $K\ge1$ may be the same sample-independent target-cost bound used
by the reconstruction. In particular,
$D_{\rm KL}(Q_{S,N,t,c}\Vert\mu_{N,t})\le NK/t^2$.
For the loss
\[
 \ell_w(x)=\mathbf1\{f(x)z_w(x)\le1/4\},
\]
the empirical loss is zero under this posterior. The following finite
reference family makes the resulting risk bound simultaneous in width
and scale; an infinite mixture over widths or exact scales is unnecessary.

\begin{lemma}[A finite Gaussian reference family]
\label{lem:si-grid}
Fix $K\ge1$ and $J\ge K$. Set
\[
 M=\lceil J/K\rceil,\qquad
 P_M=M(d+1)+(L-2)M^2+M,
\]
\[
 R_\tau=1+\left\lfloor\frac{P_M}{2}\log(J/K)\right\rfloor,
 \qquad \tau_j=\exp(-j/P_M)
          \quad(0\le j<R_\tau).
\]
There is a sample-independent prior on the disjoint union of these
$MR_\tau$ Gaussian models such that, whenever $NK/t^2\le J$, the
original posterior, embedded in one component, has relative entropy at
most
\begin{equation}\label{eq:si-grid-kl}
 NK/t^2+1+\log(MR_\tau).
\end{equation}
\end{lemma}
\begin{proof}
The hypothesis implies $N\le M$ and $t^2\ge K/J$.
Choose $j=\lfloor P_M\log(1/t)\rfloor$. Then
$0\le j<R_\tau$, $\tau_j\ge t$, and
$\log(\tau_j/t)\le1/P_M$.
At a fixed width all Gaussian covariance matrices are multiplied by
the same scalar $t^2$. Their density ratio is therefore
\[
 \log\frac{d\mu_{N,t}}{d\mu_{N,\tau_j}}(w)
 =P_N\log(\tau_j/t)
   -\frac N2(t^{-2}-\tau_j^{-2})\kappa_N(w)
 \le P_N\log(\tau_j/t)\le1.
\]
The nonnegative weighted energy $\kappa_N$ is the one defined in the
main text. Combining this pointwise inequality with
\eqref{eq:si-raw-density} bounds relative entropy with respect to
$\mu_{N,\tau_j}$ by $NK/t^2+1$.
Choose the component $(N,j)$ uniformly from the fixed finite family,
then draw from its Gaussian law. Embedding the posterior in the chosen
component adds exactly $\log(MR_\tau)$ and proves the claim.
\end{proof}

We next state the statistical interface needed from the reconstruction.
Let $\psi$ be the ramp loss used in the posterior comparison, and write
\[
 R_\psi(\Gamma)=\mathbb E_{h\sim\Gamma,X}
                         \psi(f(X)T(h(X))),\qquad
 \widehat R_{\psi,S}(\Gamma)=\frac1n\sum_{i=1}^n
                 \mathbb E_{h\sim\Gamma}\psi(f(X_i)T(h(X_i))).
\]
The readout-retention and pointwise coupling estimates give the two
inequalities in the next theorem. Thus its hypotheses can be checked
without imposing robustness on any decoded predictor.

\begin{theorem}[A simultaneous width and scale bound]
\label{thm:si-envelope}
Suppose \eqref{eq:si-raw-density} holds. Fix nonnegative $\alpha,\beta$
and a message threshold $\Delta$. Suppose there are sample-independent
constants
\[
 A\ge\max\{1,K,\Delta\},\qquad B\ge0,
\]
with the following property. For every fixed $0<a\le1$ there is a finite
class $\mathcal H_a$ satisfying
\begin{equation}\label{eq:si-class}
 \log|\mathcal H_a|\le A+Ba^4.
\end{equation}
Whenever $t\le a$ and $NK/t^2>\Delta$, the reconstruction supplies a law
$\Gamma$ on $\mathcal H_a$ such that
\begin{equation}\label{eq:si-comparison}
 \widehat R_{\psi,S}(\Gamma)\le\alpha,
 \qquad
 R(Q_{S,N,t,c})\le\beta+\alpha+R_\psi(\Gamma).
\end{equation}

Set $J=A+B$, use $M,R_\tau$ from Lemma~\ref{lem:si-grid}, and define
\[
 j_a=\left\lceil\frac12\log_2\max\{1,B/A\}\right\rceil,
 \qquad R_a=j_a+1,
 \qquad
 \Lambda=1+\log\frac{2MR_\tau R_a}{\delta}.
\]
With sample probability at least $1-\delta$, simultaneously for every
$N\ge C$ and $0<t\le1$,
\begin{equation}\label{eq:si-envelope}
 R(Q_{S,N,t,c})
 \le\beta+3\alpha+
   \frac8n\left[\min\{NK/t^2,\ A+Bt^4\}+\Lambda\right].
\end{equation}
On the same event, the branch $NK/t^2\le\Delta$ obeys the sharper direct
bound
\begin{equation}\label{eq:si-small}
 R(Q_{S,N,t,c})\le\frac2n[\Delta+\Lambda].
\end{equation}
\end{theorem}
\begin{proof}
Use the finite scale list $a_j=2^{-j/2}$, $0\le j<R_a$.
For every $0<t\le1$, one of these upper scales satisfies $t\le a_j$ and
\begin{equation}\label{eq:si-scale-grid}
 A+Ba_j^4\le4(A+Bt^4).
\end{equation}
Indeed, above the smallest scale choose the nearest upper scale, whose
square is less than twice $t^2$. Below the smallest scale,
$Ba_{j_a}^4\le A$ if $B>A$, while if $B\le A$ the single scale $a=1$
already costs at most $2A$.

Construct one prior as follows. Give weight $1/2$ to the Gaussian
reference family of Lemma~\ref{lem:si-grid}. Give the remaining weight
$1/2$ to a uniform choice of $j$, followed by the uniform law on
$\mathcal H_{a_j}$. Keep the component tag as part of the hypothesis.
For Gaussian components use the test-margin indicator loss; for finite
code components use the ramp loss. These are a single $[0,1]$-valued loss
on the tagged space. Apply Lemma~\ref{lem:si-risk} to this prior.

Put $b=NK/t^2$ and $F=A+Bt^4\le J$.
If $b\le F$, then $b\le J$, so the Gaussian family applies. The
posterior has zero empirical indicator loss, and its relative-entropy
charge, including the mixture weight and confidence term, is at most
$b+\Lambda$. Thus its risk is at most $2(b+\Lambda)/n$.

If $b>F$, then $b>A\ge\Delta$, so the reconstruction applies with the
scale selected in \eqref{eq:si-scale-grid}. Its relative-entropy charge
is at most $4F+\Lambda$. Lemma~\ref{lem:si-risk} and
\eqref{eq:si-comparison} give risk at most
$\beta+3\alpha+2(4F+\Lambda)/n$.
These two cases imply \eqref{eq:si-envelope}.
Finally, $b\le\Delta$ implies $b\le A\le J$, so the direct Gaussian
calculation gives \eqref{eq:si-small} on this same event.
\end{proof}

In particular, taking $\beta=\varepsilon/4$ and
$\alpha=\varepsilon/12$, a sufficient sample condition at any given
width and scale is
\begin{equation}\label{eq:si-sample-envelope}
 n\ge\frac{16}{\varepsilon}
       \left[\min\{NK/t^2,\ A+Bt^4\}+\Lambda\right].
\end{equation}
The event in Theorem~\ref{thm:si-envelope} is common to all such pairs.
For polynomially bounded $A,B,K,d,L$, the additional term $\Lambda$
is logarithmic in these resources and in $1/\delta$.

\begin{corollary}[Interpolation for stronger Gaussian scales]
\label{cor:si-gamma}
Under the hypotheses of Theorem~\ref{thm:si-envelope}, let $\gamma\ge1$.
On its same sample event, it suffices, simultaneously for every $N\ge C$
and $0<t^2\le N^{1-\gamma}$, to require
\begin{equation}\label{eq:si-generic-gamma}
 n\ge\frac{16}{\varepsilon}
 \left[A+
   K^{(2\gamma-2)/(3\gamma-2)}
   B^{\gamma/(3\gamma-2)}+\Lambda\right].
\end{equation}
The convention for $B=0$ is that its displayed contribution is zero.
\end{corollary}
\begin{proof}
If $B\le K$, then $A+Bt^4\le A+B$ and
$B\le K^{(2\gamma-2)/(3\gamma-2)}B^{\gamma/(3\gamma-2)}$;
the result follows directly. Suppose $B>K$ and set
\[
 M_\gamma=(B/K)^{1/(3\gamma-2)},
 \qquad u_0=M_\gamma^{1-\gamma},
 \qquad u=t^2.
\]
If $u\le u_0$, the code term is at most $A+Bu_0^2$.
If $u>u_0$ and $\gamma>1$, the restriction
$u\le N^{1-\gamma}$ implies $N<M_\gamma$, and therefore
$NK/u\le KM_\gamma/u_0=KM_\gamma^\gamma$.
These two terms agree because
\[
 Bu_0^2=KM_\gamma^\gamma
       =K^{(2\gamma-2)/(3\gamma-2)}B^{\gamma/(3\gamma-2)}.
\]
For $\gamma=1$, $u_0=1$ and only the first case is needed.
Thus the minimum in \eqref{eq:si-sample-envelope} is at most the bracket
in \eqref{eq:si-generic-gamma}, proving the claim.
\end{proof}

\paragraph{Application to the entropy-metric code.}
The new message has log-cardinality
\[
 D_0=[(d+1)q+2]\log2+
       q^2\log(3+8K_aq/e),\qquad
 K_a=2K+\log(1/\eta),\quad q=\lceil K_a/e\rceil.
\]
Take $\Delta=\max\{K,2LD_0\}$. The large-branch condition in
Theorem~\ref{thm:si-envelope} is exactly the condition needed to absorb
the accumulated messages under the conservative cutoff
$N/t^2>\max\{1,2LD_0/K\}$.
The complementary branch costs at most $\Delta=O(K+LD_0)$ by
\eqref{eq:si-small}, instead of being rounded as a dense network of
width $LD_0/K$.

For a fixed upper scale $a$, the variance padding can be chosen as
$\lambda=\min\{1,e/(c_0a^2)\}$, where $c_0$ is a sufficiently large
absolute constant. A separate uniform source sample of size
$\widetilde O(1+\lambda^{-2})=
 \widetilde O(1+a^4/e^2)$ estimates the total variance to accuracy
$\lambda$. The extra physical Gaussian variance is $O(a^2\lambda)=O(e)$.
Together with the active width $\widetilde O(K/e)$ and the quadratic
message, the construction therefore gives constants in
\eqref{eq:si-class} with
\begin{equation}\label{eq:si-AB}
 A=\widetilde O(LK^2/e^2),\qquad
 B=\widetilde O(LK/e^3),
\end{equation}
where $A$ is enlarged to include $\Delta$ and the fixed lower-order terms.
These are bounds on pre-sample constants, not on sample-dependent
descriptions. In particular, the direct small branch is absorbed by
the first term in \eqref{eq:si-AB}.

For fixed $d,c,\varepsilon$, put
$F_0=1+\log^2(2+C)$ and $G=\log(2+CL)$.
The target-cylinder estimates give
$K=O(CLF_0)$ and $e^{-1}=O(CL^3F_0)$.
Consequently $q=O(C^2L^4F_0^2)$,
$D_0=O(q^2G)$, and the total source width at upper scale $a$ is
\[
 m=O\left([C^2L^4F_0^2+a^4C^2L^6F_0^2]G\right).
\]
The scalar count is $O(L[mq+q^2])$, up to lower-order first-layer
terms. Polynomial coefficient bounds and inverse meshes contribute one
further factor $G$, uniformly in $a\le1$: in particular
$\lambda^{-1}\le c_0/e$. Thus, after increasing a constant depending
only on $d,c,\varepsilon$, the choices in \eqref{eq:si-AB} satisfy
\begin{equation}\label{eq:si-resources}
 A=O(C^4L^9F_0^4G^2),\qquad
 B=O(C^4L^{11}F_0^4G^2).
\end{equation}
Moreover $\Lambda=O(G+\log(1/\delta))$.
The width and scale envelope becomes
\[
 \widetilde O\min\{CLN/t^2,\ C^4L^9+t^4C^4L^{11}\}.
\]
Corollary~\ref{cor:si-gamma} gives the sufficient sample order
\begin{equation}\label{eq:si-new-powers}
 \widetilde O\left(
 C^4L^9+
 C^{(6\gamma-2)/(3\gamma-2)}
 L^{(13\gamma-2)/(3\gamma-2)}\right).
\end{equation}
An explicit common logarithmic factor in \eqref{eq:si-new-powers} is
$F_0^4G^2$, with an additional additive $O(\log(1/\delta))$ term
before the accuracy-dependent constant. Indeed, the interpolated term
has $F_0$ exponent $(6\gamma-2)/(3\gamma-2)\le4$ and $G$ exponent
$2\gamma/(3\gamma-2)\le2$ for every $\gamma\ge1$.
At $\gamma=1$ this is $\widetilde O(C^4L^{11})$.
For every $\gamma\ge1$ the second term's width exponent is at most four;
its depth exponent is at most nine exactly when $\gamma\ge8/7$.
Thus $\widetilde O(C^4L^9)$ suffices throughout the regime
$0<t^2\le N^{-1/7}$. The existing stronger-scale theorem remains
sharper at $0<t^2\le N^{-1}$, where it gives
$\widetilde O(C^4L^7)$.

The direct posterior argument uses only the target-cylinder denominator
and zero empirical margin loss. It therefore also applies to the optional
global clean-margin conditioning whenever the same cylinder is retained.
